\subsection{Polar decomposition}

In this section, we consider complex Hilbert spaces.

Let \(E_{1}\) and \(E_{2}\) be Hilbert spaces and \(u \in \mathcal{L}(E_{1}; E_{2})\) . The endomorphism \(u^{*}u\) of \(E_{1}\) is positive (prop. 8), hence one can form the positive element \((u^{*}u)^{1/2}\) of \(\mathcal{L}(E_{1})\) .

DEFINITION 3. --- One says that \((u^{*}u)^{1/2}\) is the absolute value of u, and one denotes it by \(|u|\) .

In the case where \(E_{1} = E_{2}\) , this definition coincides with that given in remark 9 of I, p.~119.

For an element \(u\) of \(\mathcal{L}(\mathrm{E}_{1};\mathrm{E}_{2})\) , recall (EVT, V, p.~41, def. 2) that the initial subspace of \(u\) is the closed subspace \(\mathrm{Ker}(u)^{\circ}\) of \(\mathrm{E}_{1}\) and the final subspace of \(u\) is the closed subspace \(\overline{\mathrm{Im}(u)}\) of \(\mathrm{E}_{2}\) .

PROPOSITION 10. --- Let \(\mathrm{E}_1\) and \(\mathrm{E}_2\) be complex Hilbert spaces and \(u \in \mathcal{L}(\mathrm{E}_1; \mathrm{E}_2)\) .

\begin{enumerate}
\def\labelenumi{\alph{enumi})}
\item
  The initial subspace and the final subspace of \(|u|\) are both equal to the initial subspace of \(u\) and we have \(\| |u| \| = \| u\|\) ;
\item
  There exists a unique partial isometry \(j\) of \(\mathrm{E}_1\) into \(\mathrm{E}_2\) such that \(\operatorname{Ker}(j) = \operatorname{Ker}(u)\) and \(u = j|u|\) ;
\item
  The initial (resp. final) subspace of j is equal to that of u;
\item
  Let \(u_{1}\) be a positive element of \(\mathcal{L}(\mathrm{E}_{1})\) and \(j_{1}\) be a partial isometry of \(\mathcal{L}(\mathrm{E}_{1};\mathrm{E}_{2})\) such that \(u = j_{1}u_{1}\) and \(\operatorname{Ker}(j_{1}) = \operatorname{Ker}(u_{1})\) . Then \(u_{1} = |u|\) and \(j_{1} = j\) .
\end{enumerate}

For every \(x\in \mathrm{E}_1\) , one has

\[
\| u (x) \| ^ {2} = \langle x | (u ^ {*} u) (x) \rangle = \langle x | | u | ^ {2} (x) \rangle = \| | u | (x) \| ^ {2}.\tag{2}
\]

This proves that \(\operatorname{Ker}(u) = \operatorname{Ker}(|u|)\) and \(\| |u| \| = \|u\|\) . Since \(|u|\) is Hermitian, the closure of the image of \(|u|\) is the orthogonal complement of its kernel (EVT, V, p.~41, prop. 4), that is, the initial space of \(u\) , whence \(a\) ).

Formula (2) implies that there exists an isometric map \(v\) of \(\operatorname{Im}(|u|)\) on \(\operatorname{Im}(u)\) such that \(v(|u|(x)) = u(x)\) for all \(x \in \mathrm{E}_1\) . Let \(j\) be the unique element of \(\mathcal{L}(\mathrm{E}_1; \mathrm{E}_2)\) which extends \(v\) and vanishes on \(\operatorname{Im}(|u|)^{\circ} = \operatorname{Ker}(|u|) = \operatorname{Ker}(u)\) . Then \(j\) has the properties of \(b\) ). The uniqueness of \(j\) follows from the decomposition \(\mathrm{E} = \operatorname{Ker}(u) \oplus \overline{\operatorname{Im}(|u|)}\) .

The initial subspace of \(j\) is \(\mathrm{Ker}(j)^{\circ} = \mathrm{Ker}(u)^{\circ}\) , the initial space of \(u\) . Its final subspace is \(\overline{j(\mathrm{Ker}(u)^{\circ})} = \overline{j(\mathrm{Im}(|u|))} = \overline{\mathrm{Im}(u)}\) , the final subspace of \(u\) . This proves \(c\) ).

Now let \(u_1\) and \(j_1\) be as in \(d\) ). We have \(u^*u = u_1j_1^*j_1u_1\) . The map \(j_1^*j_1\) is the orthogonal projector with kernel \(\text{Ker}(j_1) = \text{Ker}(u_1)\) (EVT, V, p.~41, prop. 5 (ii)) and hence with image \(\overline{\text{Im}(u_1)}\) . Therefore \(u^*u = u_1^2\) and consequently \(u_1 = (u^*u)^{1/2} = |u|\) (I, p.~118, prop. 16). The uniqueness assertion of \(b\) ) finally implies that \(j_1 = j\) .

DEFINITION 4. --- Let \(E_{1}\) and \(E_{2}\) be complex Hilbert spaces and \(u \in \mathcal{L}(E_{1}; E_{2})\) . The pair \((j, |u|)\) , where j is the unique partial isometric map from \(E_{1}\) into \(E_{2}\) such that \(u = j|u|\) and \(\operatorname{Ker}(j) = \operatorname{Ker}(u)\) , is called the polar decomposition of u.

PROPOSITION 11. — Let \(\mathrm{E}_1\) and \(\mathrm{E}_2\) be complex Hilbert spaces and \(u \in \mathcal{L}(\mathrm{E}_1; \mathrm{E}_2)\) . Let \((j, |u|)\) be the polar decomposition of \(u\) .

\begin{enumerate}
\def\labelenumi{\alph{enumi})}
\item
  We have \(|u| = j^{*}u = u^{*}j\) ;
\item
  We have \(|u^{*}| = j u^{*} = u j^{*}\) ;
\item
  The polar decomposition of \(u^{*}\) is \((j^{*}, |u^{*}|)\) .
\end{enumerate}

Let I = \(\operatorname{Ker}(u)^{\circ}\) and \(F = \overline{\operatorname{Im}(u)}\) be the initial subspace and the final subspace of u; moreover, we have I = \(\operatorname{Ker}(|u|)^{\circ} = \overline{\operatorname{Im}(|u|)}\) (prop. 10, a)). The map \(j^{*}j\) is the orthogonal projector of \(E_{1}\) onto I (loc. cit. and

EVT, V, p.~41, prop. 5 (ii)). We therefore have \(j^{*}u = j^{*}j|u| = |u|\) , and then \(u^{*}j = (j^{*}u)^{*} = |u|^{*} = |u|\) , whence \(a\) ).

Similarly, we compute \(u^{*} = |u|j^{*} = (j^{*}j|u|)j^{*} = j^{*}(j|u|j^{*})\) . The endomorphism \(j|u|j^{*}\) of \(E_{2}\) is positive. The linear map \(j^{*}\) is partially isometric with initial space F and final space I (EVT, V, p.~41, prop. 5), and the linear maps from I to F (resp. from F to I) deduced from j and \(j^{*}\) by passing to subspaces are mutually inverse isomorphisms (loc. cit.). We therefore have \(\operatorname{Ker}(j|u|j^{*}) = \operatorname{Ker}(|u|j^{*}) = \operatorname{Ker}(j^{*})\) , since the image of \(j^{*}\) is contained in \(\operatorname{Ker}(u)^{\circ} = \operatorname{Ker}(|u|)^{\circ}\) . By prop. 10, d), the pair \((j^{*}, j|u|j^{*})\) is the polar decomposition of \(u^{*}\) . This proves c), and assertion b) then follows from assertion a) applied to \(u^{*}\) .

COROLLARY. --- Let \(E_{1}\) and \(E_{2}\) be complex Hilbert spaces and let \(u \in \mathcal{L}(E_{1}; E_{2})\) . We have \(\operatorname{Im}(u) = \operatorname{Im}(|u^{*}|)\) .

Let \((j, |u|)\) be the polar decomposition of \(u\) . We have \(|u^{*}| = j|u|j^{*} = uj^{*}\) by the preceding proposition. By EVT, V, p.~41, prop. 5, the map \(j^{*}\) is partially isometric. Its final space is \(\operatorname{Ker}(j)^{\circ} = \operatorname{Ker}(u)^{\circ}\) (prop. 10, c)). The assertion follows.

PROPOSITION 12. — Let \(E_{1}\) and \(E_{2}\) be complex Hilbert spaces and \(u \in \mathcal{L}(E_{1}; E_{2})\) . Let \((j, |u|)\) be the polar decomposition of u. For u to be bijective, it is necessary and sufficient that \(|u|\) be invertible in \(\mathcal{L}(E_{1})\) and that j be an isomorphism from \(E_{1}\) onto \(E_{2}\) .

The condition is sufficient. Conversely, if u is bijective, then \(u^{*}u\) is invertible in \(\mathcal{L}(\mathrm{E}_{1})\) , and \(|u| = (u^{*}u)^{1/2}\) is also bijective since its spectrum is contained in \(R_{+}^{*}\) . Moreover, \(\operatorname{Ker}(j) = \operatorname{Ker}(u) = \{0\}\) and \(\operatorname{Im}(j) = \operatorname{Im}(u) = F\) , so j maps \(E_{1}\) isometrically onto \(E_{2}\) .

PROPOSITION 13. — Let E be a complex Hilbert space and let u be an endomorphism of E. The following conditions are equivalent:

\begin{enumerate}
\def\labelenumi{(\roman{enumi})}
\item
  The endomorphism u is normal;
\item
  There exists a unitary element v of \(\mathcal{L}(\mathrm{E})\) , commuting with \(|u|\) , such that \(u = v|u|\) .
\end{enumerate}

Let \((j,|u|)\) be the polar decomposition of \(u\) . Suppose that \(u\) is normal. Then we have \(|u^{*}| = (uu^{*})^{1/2} = (u^{*}u)^{1/2} = |u|\) . Prop. 11 then implies \(|u|j = |u^{*}|j = ju^{*}j = j|u|\) . Moreover, \(j\) stabilizes the supplementary orthogonal subspaces \(\operatorname{Ker}(|u|) = \operatorname{Ker}(j)\) and \(\overline{\operatorname{Im}(|u|)} = \overline{\operatorname{Im}(j)}\) (prop. 10). Let \(v\) be the element of \(\mathcal{L}(\mathrm{E})\) which coincides with \(j\) on \(\operatorname{Ker}(u)^{\circ}\) and with the identity map on \(\operatorname{Ker}(u)\) . Since \(j\) induces an isometry of \(\operatorname{Ker}(u)^{\circ}\) onto \(\overline{\operatorname{Im}(u)} = \operatorname{Ker}(u^{*})^{\circ} = \operatorname{Ker}(u)^{\circ}\) (since \(u\) is normal), the endomorphism \(v\) is unitary; moreover, it commutes with \(|u|\) since \(j|u| = |u|j\) , and we have \(u = v|u|\) .

Conversely, let v be a unitary element of \(\mathcal{L}(\mathrm{E})\) , commuting with \(|u|\) , such that \(u = v|u|\) . We have \(uu^{*} = v|u|^{2}v^{*} = |u|^{2}vv^{*} = |u|^{2} = u^{*}u\) , hence u is normal.

\pandocbounded{\includegraphics[keepaspectratio]{images/part001_8ac1eb40.jpg}}

Let E be a complex Hilbert space and \(u \in \mathcal{L}(\mathrm{E})\) . Let \((j, |u|)\) be the polar decomposition of u. It is possible that j commutes with \(|u|\) without u being normal (exercise 11 of I, p.~189).

